\documentclass[11pt,a4paper]{article}
\usepackage[T1]{fontenc}
\usepackage[utf8]{inputenc}
\usepackage[english]{babel}
\usepackage{amsmath,amssymb,amsthm,mathtools}
\usepackage[margin=27mm]{geometry}
\usepackage{microtype}
\usepackage[colorlinks=true,linkcolor=blue,citecolor=blue,urlcolor=blue]{hyperref}
\newtheorem{theorem}{Theorem}[section]
\newtheorem{proposition}[theorem]{Proposition}
\newtheorem{lemma}[theorem]{Lemma}
\newtheorem{candidate}[theorem]{Candidate theorem}
\theoremstyle{definition}
\newtheorem{definition}[theorem]{Definition}
\theoremstyle{remark}
\newtheorem{remark}[theorem]{Remark}
\newcommand{\R}{\mathbb R}
\newcommand{\C}{\mathbb C}
\newcommand{\Sph}{\mathbb S^2}
\newcommand{\vol}{\operatorname{Vol}_3}
\newcommand{\area}{\mathcal H^2}
\newcommand{\F}{\mathcal F}
\newcommand{\supp}{\operatorname{supp}}
\newcommand{\diag}{\operatorname{diag}}
\newcommand{\id}{\operatorname{Id}}
\DeclareMathOperator{\Rea}{Re}
\DeclareMathOperator{\Ima}{Im}
\numberwithin{equation}{section}
\title{Covariogram Uniqueness with One Regular Body in $\R^3$\\[3pt]
\large A Candidate Proof via Area Measures and Fourier--Laplace Zeros}
\author{Andrea\\[2pt]\small A curious amateur with proof obligations}
\date{9 October 2026\\\small Working draft for mathematical review}

\begin{document}
\maketitle
\begin{abstract}
Does the covariogram of a $C^9_+$ convex body in $\R^3$ determine it among
all convex bodies, without a regularity assumption on the competitor?
This extension is explicitly identified as open in Bianchi's survey.
We present a candidate proof based on two proposed steps. First, the common
difference body yields a $C^{1,1}$ support function for the competitor;
quadratic data almost everywhere and Laplace asymptotics then recover
unordered antipodal area densities. Second, analytic allocation of the
regular body's Fourier--Laplace zeros and a positive weighted limit impose
one common orientation on these densities. Minkowski uniqueness supplies
the final implication. The main text states the argument and its dependencies;
the appendices give the detailed proposed proofs. This is a submission for
independent scrutiny, not a claim of a validated solution. A separate
supplement records partial Lean verification, which does not certify the
two central propositions or the complete uniqueness theorem.
\end{abstract}

\section{Introduction and purpose}\label{sec:introduction}
The problem addressed here was chosen because it was explicitly declared
open. Bianchi's uniqueness theorem applies when both bodies have sufficiently
regular boundaries \cite[Theorem~1.1 and Remark~5.2]{Bianchi2016}. In dimension
three the hypothesis is $C^9_+$. His survey identifies the remaining question:
whether such a body is determined among \emph{all} convex bodies
\cite[discussion following Theorem~3.25, manuscript p.~23]{Bianchi2023}.
The absence of a regularity hypothesis on the second body is the issue in
this paper.

The investigation began as an experiment in AI-assisted mathematical
exploration by an author outside professional research in convex geometry.
Its purpose is to bring the resulting candidate argument to a form that
a specialist can evaluate. The experience also motivates a broader question:
what could an expert accomplish by combining these tools with the knowledge
needed to choose productive directions and detect subtle errors?
The interaction record and computational artifacts belong to the accompanying
research materials; the present paper concentrates on the mathematics.

A convex body is a compact convex set with nonempty interior. Its
covariogram is
\begin{equation}\label{eq:covariogram}
 g_C(x)=\vol(C\cap(C+x))
       =\int_{\R^3}\mathbf1_C(y)\mathbf1_C(y-x)\,dy .
\end{equation}
Translations and central reflections preserve this function. We write
$C\in C^r_+$ when $\partial C$ is $C^r$ and has positive Gaussian curvature.
The precise proposed conclusion is the following.

\begin{candidate}\label{thm:target}
Let $K\subset\R^3$ be a convex body of class $C^9_+$, and let
$L\subset\R^3$ be an arbitrary convex body. If $g_L=g_K$, then
\[
 L=v+K\qquad\text{or}\qquad L=v-K
 \quad\text{for some }v\in\R^3 .
\]
\end{candidate}

The proposed proof has two substantive components:
Proposition~\ref{prop:geometric} transfers geometric information to the
possibly nonsmooth competitor; Proposition~\ref{prop:transport} uses only
the regular body's spectral information to make the orientation global.
Their detailed arguments are given in Appendices~\ref{app:geometry}--\ref{app:mean}.
Section~\ref{sec:completion} shows exactly how they imply the candidate
theorem. These two propositions are submitted for review as proposed
steps, rather than cited as established results.

The target is the extension identified in the cited literature. The
bibliographic investigation recorded with the research materials found
no subsequent resolution in the sources consulted as of 9 October 2026;
this is not an exhaustive priority claim. The argument does not address
an arbitrary first body, and no lower-regularity extension is asserted here.

\section{Geometric reduction for the competitor}\label{sec:geometry}
Let $h_C$ denote the positively homogeneous support function of $C$ and set
\begin{align}
 \alpha_C(u)&=-h_C(-u),&
 \beta_C(u)&=h_C(u),&
 w_C(u)&=\beta_C(u)-\alpha_C(u),\label{eq:endpoints}\\
 A_C(u,t)&=\area(C\cap\{x:x\cdot u=t\}),&&u\in\Sph .\label{eq:sections}
\end{align}
We use $\sigma=\area|_{\Sph}$, and $S_C$ denotes the surface area measure.
For a $C^2_+$ body,
\begin{equation}\label{eq:radii}
 R_C(u)=\nabla_{\Sph}^2h_C(u)+h_C(u)\id,\qquad
 f_C(u)=\det R_C(u)=\kappa_C(u)^{-1},\qquad dS_C=f_C\,d\sigma .
\end{equation}
Determinants are taken on the two-dimensional tangent plane.
For a $C^{1,1}$ support function, $R_C$ is interpreted almost everywhere.

Equal covariograms give the same difference body and volume:
\begin{equation}\label{eq:difference}
 D:=K-K=L-L,\qquad w_K=w_L=:w,\qquad \vol(K)=\vol(L).
\end{equation}
Indeed $\supp g_C=C-C$, with support understood as the closure of the
nonzero set, and $g_C(0)=\vol(C)$.
Define the entire transform
\begin{equation}\label{eq:transform}
 \F_C(z)=\int_C e^{iz\cdot x}\,dx,\qquad z\in\C^3 .
\end{equation}
The convolution identity in \eqref{eq:covariogram} yields
\begin{equation}\label{eq:product}
 \F_L(z)\F_L(-z)=\F_K(z)\F_K(-z),\qquad z\in\C^3.
\end{equation}
The scalar restrictions also satisfy
\begin{equation}\label{eq:reality}
 \F_C(-\overline\zeta u)=\overline{\F_C(\zeta u)},\qquad
 \F_C(i\tau u)>0\quad(\tau\in\R).
\end{equation}

\begin{proposition}[Geometric transfer; proposed]\label{prop:geometric}
Suppose $K\in C^2_+$ and $g_L=g_K$, with $L$ an arbitrary convex body.
Then $h_L\in C^{1,1}_{\mathrm{loc}}(\R^3\setminus\{0\})$ and
\begin{equation}\label{eq:areaL}
 dS_L=f_L\,d\sigma,\qquad f_L(u)=\det R_L(u).
\end{equation}
For almost every $u$, $R_L(u)$ and $R_L(-u)$ are positive definite and
\begin{equation}\label{eq:unordered}
 \{f_L(u),f_L(-u)\}=\{f_K(u),f_K(-u)\}.
\end{equation}
At such directions the endpoint expansions are
\begin{align}
 A_L(u,\alpha_L+s)&=a_L(u)s+o(s),&
 a_L(u)&=2\pi\sqrt{f_L(-u)},\label{eq:lower}\\
 A_L(u,\beta_L-s)&=b_L(u)s+o(s),&
 b_L(u)&=2\pi\sqrt{f_L(u)},\label{eq:upper}
\end{align}
as $s\downarrow0$.
\end{proposition}

The proposed argument first uses $0\le D^2h_L\le D^2h_D$ in the sense
of matrix-valued measures. This gives a Lipschitz map
$u\mapsto\nabla h_L(u)$ and, through the area formula, the entire measure
identity \eqref{eq:areaL}. An endpoint section calculation then gives
\eqref{eq:lower}--\eqref{eq:upper}, initially with nonnegative coefficients.
The imaginary-axis identity \eqref{eq:product} forces their product to
equal the strictly positive product for $K$. Finally, the directional
derivatives of $g_C$ at zero recover the even part of $S_C$.
Thus products and sums recover \eqref{eq:unordered}.
Appendix~\ref{app:geometry} supplies the details, including degenerate
quadratic forms and the possible singular part of $S_L$.

\section{A common orientation from the spectral data}\label{sec:spectral}
For the regular body define the continuous odd function
\begin{equation}\label{eq:gammaK}
 \gamma_K(u)=\frac{\log f_K(u)-\log f_K(-u)}{2w(u)} .
\end{equation}
The spectral input used below is Bianchi's extension of Kobayashi's
description of the zeros \cite[Theorems~4.1--4.2]{Bianchi2016}.
For $K\in C^9_+$, there are real-analytic branches
$z_m:\Sph\to\C$, for all sufficiently large integers $m$, such that
\begin{equation}\label{eq:rootsK}
 z_m(u)=\frac{(2m+1)\pi}{w(u)}+i\gamma_K(u)+O(m^{-1})
\end{equation}
uniformly in $u$. They are distinct simple zeros of
$\zeta\mapsto\F_K(\zeta u)$ with positive real part.
The remaining zeros with positive real part lie in a bounded set common
to all directions, and are finite in number in each fixed direction,
counted with multiplicity. Simplicity and the uniform bound also appear
in the proof of the cited Theorem~4.2, in particular its equation~(4.16).

If $K$ is not centrally symmetric, then $\gamma_K\not\equiv0$.
Otherwise $S_K$ is even, and Minkowski uniqueness gives a translate of
$-K$ equal to $K$.

\begin{proposition}[Transport of endpoint ratios; proposed]\label{prop:transport}
Let $K\in C^9_+$ be noncentrally symmetric and let $g_L=g_K$.
Assume that \eqref{eq:lower}--\eqref{eq:upper} hold with positive
finite coefficients on a set of full $\sigma$-measure.
Then there is a unique $c\in[-1,1]$, independent of direction, such that
\begin{equation}\label{eq:transport}
 \gamma_L(u):=\frac{\log(b_L(u)/a_L(u))}{w(u)}
             =c\gamma_K(u)\quad\text{for almost every }u.
\end{equation}
\end{proposition}

Here is the mechanism. The analytic functions
$\F_L(z_m(u)u)$ and $\F_L(-z_m(u)u)$ have identically zero product.
Analytic continuation and simplicity therefore allocate each sufficiently
large branch to one factor globally. Consequently, a sequence
$\varepsilon_m\in\{-1,1\}$ independent of $u$ describes the positive-real-part
zeros of the competitor, apart from finitely many:
\begin{equation}\label{eq:rootsL}
 \zeta_{m,L}=x_m+iy_m,\qquad
 x_m=\frac{(2m+1)\pi}{w}+O(m^{-1}),\qquad
 y_m=\varepsilon_m\gamma_K(u)+O(m^{-1}).
\end{equation}
The allocation, including real zeros, is justified in Appendix~\ref{app:allocation}.

For a fixed admissible direction write $F(\zeta)=\F_L(\zeta u)$ and define
\begin{equation}\label{eq:DH}
 D_L(\tau)=\log F(i\tau)-\log F(-i\tau),\qquad
 H_L(\tau)=D_L(2\tau)-2D_L(\tau).
\end{equation}
Endpoint Laplace asymptotics give
$H_L(\tau)\to\log(b_L/a_L)$. The finite difference cancels the linear
exponential term in a Hadamard product. Pairing the zeros as
$x_j+iy_j,-x_j+iy_j$ and using the estimates in Appendix~\ref{app:kernels}
gives
\begin{equation}\label{eq:Hkernel}
 H_L(\tau)=
 \sum_j\frac{24y_j\tau^3}
 {(x_j^2+\tau^2)(x_j^2+4\tau^2)}+o(1).
\end{equation}
After \eqref{eq:rootsL} and the rescaling $\tau=2\pi R/w(u)$,
\begin{equation}\label{eq:Hmean}
 H_L(2\pi R/w(u))=w(u)\gamma_K(u)N(R)+o(1),
\end{equation}
where
\begin{equation}\label{eq:mean}
 N(R)=\frac{12}{\pi}\sum_{m\ge m_0}
 \frac{\varepsilon_m R^3}
 {((m+\tfrac12)^2+R^2)((m+\tfrac12)^2+4R^2)} .
\end{equation}
The weights are positive and have total mass tending to one.
One admissible direction with $\gamma_K(u)\ne0$ forces $N(R)$ to have
a limit $c$. The same sequence $N$ then gives \eqref{eq:transport} in
every admissible direction. Appendix~\ref{app:mean} proves these assertions
without differentiating an infinite series.

\section{Completion of the candidate argument}\label{sec:completion}
Assume the two proposed propositions with their proofs in the appendices.
If $K$ is noncentrally symmetric, Proposition~\ref{prop:geometric} gives
the hypotheses of Proposition~\ref{prop:transport} and
\begin{equation}\label{eq:localsign}
 \gamma_L(u)\in\{\gamma_K(u),-\gamma_K(u)\}
 \quad\text{almost everywhere}.
\end{equation}
Because $\gamma_K$ is continuous and nonzero somewhere, its nonzero set
contains a nonempty open set. This set meets the full-measure set where
\eqref{eq:transport} and \eqref{eq:localsign} both hold, so $c=1$ or $c=-1$.

If $c=1$, the ordered logarithmic ratio and the unordered pair
\eqref{eq:unordered} imply $f_L(u)=f_K(u)$ almost everywhere.
If $c=-1$, they imply $f_L(u)=f_K(-u)$ almost everywhere.
Where the two values of $f_K$ coincide there is no choice to make.
Identity \eqref{eq:areaL} is an identity of the complete measures, hence
$S_L=S_K$ or $S_L=S_{-K}$. The uniqueness part of Minkowski's theorem
\cite{Schneider} gives the desired translation or central reflection.

If $K$ is centrally symmetric, \eqref{eq:difference} gives
$\vol(D)=8\vol(K)=8\vol(L)$. Equality therefore holds in the
Brunn--Minkowski inequality for $L+(-L)$. Its equality case implies that
$L$ and $-L$ are homothetic; their equal volumes make them translates.
Thus $L$ is centrally symmetric, and the common difference body determines
both bodies up to translation. This completes the proposed deduction of
Candidate Theorem~\ref{thm:target}.

\section{Scope of verification and accompanying materials}\label{sec:verification}
The mathematical claims requiring scrutiny are the geometric transfer in
Appendix~\ref{app:geometry} and the spectral transport through
Appendices~\ref{app:allocation}--\ref{app:mean}. Particular attention is
needed at the area-formula step, the degenerate endpoint calculation,
the allocation of multiplicities, and the passage from the canonical
product to the directional limit. No regularity of $L$ is assumed in
these arguments beyond what is explicitly derived.

Supplement S1 records 33 supporting lemmas and conditional deductions
checked in Lean 4.24.0. It includes the exact statements, dependencies,
and limits of that verification. Neither Proposition~\ref{prop:geometric}
nor Proposition~\ref{prop:transport}, nor the complete candidate theorem,
has been formally certified. Supplement S2 describes numerical diagnostics.
All mathematical steps needed to assess the candidate are written here;
the reader need not inspect Lean code or the interaction history.

\paragraph{AI assistance.}
AI tools assisted substantially in exploration, drafting, computation and
partial formalization. The source files, verification records and available
research documentation are separate accompanying materials. Their purpose
is to make the work inspectable; responsibility for the submitted claims
remains with the human author.

\appendix
\section{Geometric transfer and endpoint sections}\label{app:geometry}
This appendix gives the proposed proof of Proposition~\ref{prop:geometric}.
The external tools are the differentiability and area theorems for Lipschitz
maps, standard convex duality, and the classical area-measure results in
convex geometry; see \cite{EvansGariepy,Schneider}. Their applications to
the present competitor are made explicit below.

\subsection{The support function and the complete area measure}
Since $K\in C^2_+$, its homogeneous support function is $C^2$ away from
zero. The same is true of $h_D=h_K+h_{-K}$. Convexity gives positive
semidefinite distributional Hessians and
\begin{equation}\label{eq:hessian-domination}
 0\le D^2h_L\le D^2h_D,\qquad h_D=h_L+h_{-L}.
\end{equation}
On a relatively compact ball away from zero, $D^2h_D$ is a bounded
continuous matrix field. Domination of the quadratic forms in
\eqref{eq:hessian-domination} first bounds the diagonal measures of
$D^2h_L$ by a multiple of Lebesgue measure; polarization then bounds
the off-diagonal entries as well. Thus the weak Hessian of $h_L$ is
locally $L^\infty$. Equivalently, mollification on a slightly larger
ball gives uniformly bounded positive semidefinite Hessians and uniformly
Lipschitz gradients. Passing to the limit shows that the continuous
convex representative $h_L$ is $C^{1,1}$.

Every support face of $L$ is consequently a singleton. Write
$p_L(u)=\nabla h_L(u)$. This map is Lipschitz on the compact sphere
and maps it onto $\partial L$, because every boundary point admits
a supporting normal. Almost everywhere,
\[
 Dp_L(u)|_{u^\perp}=R_L(u)
   =\nabla_{\Sph}^2h_L(u)+h_L(u)\id\succeq0 .
\]
The tangential two-dimensional Jacobian is $\det R_L(u)$.

Let $\partial L_{\rm reg}$ be the set of boundary points with a unique
outer unit normal $\nu_L(x)$. The complement has $\area$-measure zero.
For a Borel set $E\subset\Sph$, the multiplicity of the map $p_L|_E$
at a point $x\in\partial L_{\rm reg}$ is exactly
$\mathbf1_E(\nu_L(x))$: every preimage is a supporting normal at $x$,
and the normal is unique. The area formula therefore gives
\begin{align}
 \int_E\det R_L(u)\,d\sigma(u)
 &=\int_{\partial L}\#\{u\in E:p_L(u)=x\}\,d\area(x)\notag\\
 &=\int_{\partial L_{\rm reg}}\mathbf1_E(\nu_L(x))\,d\area(x)
 =S_L(E).\label{eq:area-formula}
\end{align}
The integral over the nonregular boundary is zero, even if its
pointwise multiplicity is infinite. This proves the identity of the
whole measure in \eqref{eq:areaL}, and not merely a formula for an
absolutely continuous part.

\subsection{The section coefficient at a quadratic direction}\label{app:sections}
Rademacher's theorem applied to the Lipschitz gradient, in a finite
collection of local charts, gives a second-order expansion of $h_L$
at almost every direction. Indeed, integration of the first-order
expansion of the gradient along a segment gives the quadratic Taylor
formula. Intersect this full-measure set with its antipodal image.

Fix a direction in this set. Rotate it to $e_3$, and translate its
support point to zero. Write a point of $\R^3$ as $(y,z)$, with
$y\in\R^2$. Then $\beta_L(e_3)=0$ and
\begin{equation}\label{eq:G-expansion}
 G(v):=h_L(v,1)=\tfrac12v^TRv+o(|v|^2),\qquad
 R=R_L(e_3)\succeq0 .
\end{equation}
We claim that
\begin{equation}\label{eq:section-coefficient}
 \lim_{s\downarrow0}\frac{A_L(e_3,-s)}s
   =2\pi\sqrt{\det R},
\end{equation}
including the case $\det R=0$.

The support face at height zero is the singleton $\{0\}$. Compactness
therefore implies that all points of the sections at height $-s$
tend to zero as $s\downarrow0$. Choose $C>0$ so that
$G(v)\le C|v|^2$ near zero. The support inequality for $(y,-s)\in L$ is
$v\cdot y-s\le G(v)$. For small sections we may choose $v=y/(2C)$,
which gives
\begin{equation}\label{eq:section-bound}
 |y|^2\le4Cs.
\end{equation}
Thus the rescaled planar sets
$E_s=\{Y:(\sqrt{s}Y,-s)\in L\}$ lie in a common disk.

Suppose first that $R$ is singular. If $s_j\downarrow0$ and
$Y_j\in E_{s_j}$ converge to $Y$, substitute $v=\sqrt{s_j}V$
in the support inequality and divide by $s_j$ to obtain
\begin{equation}\label{eq:limit-support}
 V\cdot Y-1\le\tfrac12 V^TRV\qquad(V\in\R^2).
\end{equation}
For $V$ an arbitrary real multiple of a vector in $\ker R$, this forces
$Y\perp\ker R$, that is, $Y\in\operatorname{im}R$.
Every limit point of the bounded sets $E_s$ lies in a subspace of
dimension at most one. For each $\delta>0$, these sets must eventually
lie in the $\delta$-neighborhood of this subspace within the common
disk: otherwise a convergent subsequence would contradict
\eqref{eq:limit-support}. The areas of these neighborhoods tend to zero
as $\delta\downarrow0$. Since $A_L(e_3,-s)/s=\area(E_s)$, this proves
\eqref{eq:section-coefficient} in the singular case.

Suppose now that $R$ is positive definite. The point $0\in\partial L$
has only one normal. To see this, an additional normal, combined with
$e_3$ in the convex normal cone, would give directions $(v,1)$ with
arbitrarily small nonzero $v$ for which $G(v)=0$. This contradicts the
positive quadratic lower bound in \eqref{eq:G-expansion}.
An opposite vertical normal is excluded by the nonempty interior of $L$.

Let $P$ be the projection of $L$ onto $\R^2$. The origin is in
$\operatorname{int}P$: a supporting line to $P$ at zero would otherwise
give another, horizontal, normal to $L$ at zero. The upper boundary
over $P$ is $z=-\phi(y)$, where $\phi$ is convex and finite near zero,
$\phi(0)=0$, and $\partial\phi(0)=\{0\}$. Extend $\phi$ by $+\infty$
outside $P$. It is a closed proper convex function and
\[
 G(v)=\sup_{y\in P}\{v\cdot y-\phi(y)\}.
\]
Subgradients of $\phi$ near zero are locally bounded. By the closedness
of their graph and $\partial\phi(0)=\{0\}$, every choice
$v_y\in\partial\phi(y)$ tends to zero as $y\to0$.

For $0<\epsilon<\lambda_{\min}(R)$, \eqref{eq:G-expansion} gives
\begin{equation}\label{eq:G-sandwich}
 \tfrac12v^T(R-\epsilon\id)v\le G(v)
 \le\tfrac12v^T(R+\epsilon\id)v
\end{equation}
for all sufficiently small $v$. The conjugate inequality
$\phi(y)\ge v\cdot y-G(v)$, evaluated at
$v=(R+\epsilon\id)^{-1}y$, yields the lower bound below.
For the upper bound take $v_y\in\partial\phi(y)$ and use
$\phi(y)=v_y\cdot y-G(v_y)$ and the lower bound in
\eqref{eq:G-sandwich}. Completing the square gives
\begin{equation}\label{eq:phi-sandwich}
 \tfrac12y^T(R+\epsilon\id)^{-1}y
 \le\phi(y)\le
 \tfrac12y^T(R-\epsilon\id)^{-1}y
\end{equation}
near zero. These arguments justify the use of local quadratic estimates
without assuming a globally quadratic support function.

The lower endpoint $\ell(y)$ of the vertical fiber of $L$ is a finite
convex function on $\operatorname{int}P$, hence continuous there.
The fiber over zero has positive length, since
$\operatorname{int}P$ is the projection of $\operatorname{int}L$.
Thus $\ell(y)<-d$ near zero for some $d>0$.
Together with the localization of the sections, this means that for
small $s$ the section at height $-s$ is exactly the local sublevel set
$\{\phi\le s\}$; the lower endpoint of the fiber does not restrict it.

The ellipses in \eqref{eq:phi-sandwich} now give
\[
 2\pi\sqrt{\det(R-\epsilon\id)}
 \le \liminf_{s\downarrow0}\frac{A_L(e_3,-s)}s
 \le \limsup_{s\downarrow0}\frac{A_L(e_3,-s)}s
 \le2\pi\sqrt{\det(R+\epsilon\id)} .
\]
Here we used the area formula
$\area\{y:y^TT^{-1}y/2\le s\}=2\pi s\sqrt{\det T}$
for a positive definite $2\times2$ matrix $T$.
Letting $\epsilon\downarrow0$ proves \eqref{eq:section-coefficient}.
Applying the calculation also at $-u$ establishes
\eqref{eq:lower}--\eqref{eq:upper}, at this stage with coefficients
which are finite and nonnegative.

\subsection{Laplace limits, including a zero coefficient}\label{app:laplace}
We record explicitly the endpoint fact needed before positivity is known.
If $A$ is bounded on $[\alpha,\beta]$ and
$A(\alpha+s)=as+o(s)$ with $a\ge0$, then
\begin{equation}\label{eq:abel}
 \tau^2 e^{\alpha\tau}
 \int_\alpha^\beta A(t)e^{-\tau t}\,dt\longrightarrow a .
\end{equation}
Choose $\delta>0$ so that $|A(\alpha+s)|\le Cs$ for $0<s<\delta$.
On this interval the substitution $r=\tau s$ gives an integrand
$\tau A(\alpha+r/\tau)e^{-r}$ bounded by $Cr e^{-r}$ and converging
pointwise to $ar e^{-r}$. Dominated convergence and
$\int_0^\infty re^{-r}\,dr=1$ give the limit $a$. The remaining
integral has modulus at most
$\tau^2(\beta-\alpha)\|A\|_\infty e^{-\tau\delta}$ and tends to zero.
This also proves \eqref{eq:abel} when $a=0$.
The upper endpoint has the analogous formula with $e^{-\beta\tau}$
and $e^{\tau t}$.

Sections of a fixed compact body are bounded. Fubini's theorem gives
$\F_L(\zeta u)=\int_{\alpha_L}^{\beta_L}A_L(u,t)e^{i\zeta t}\,dt$.
Consequently
\begin{equation}\label{eq:scaled-laplace}
 \tau^2e^{\alpha_L\tau}\F_L(i\tau u)\to a_L,\qquad
 \tau^2e^{-\beta_L\tau}\F_L(-i\tau u)\to b_L.
\end{equation}
The regular body's endpoint coefficients, obtained by the same positive
quadratic calculation, are
$a_K=2\pi\sqrt{f_K(-u)}$ and $b_K=2\pi\sqrt{f_K(u)}$.
Using \eqref{eq:product} and $w_L=w_K$ therefore yields
\begin{equation}\label{eq:slopes-product}
 a_Lb_L
 =\lim_{\tau\to\infty}\tau^4e^{-w\tau}
                  \F_L(i\tau u)\F_L(-i\tau u)
 =a_Kb_K>0 .
\end{equation}
Both coefficients for $L$ are positive. Its positive semidefinite
matrices $R_L(\pm u)$ thus have positive determinant and are positive
definite. Squaring the equality of products gives
\begin{equation}\label{eq:density-product}
 f_L(u)f_L(-u)=f_K(u)f_K(-u)\quad\text{almost everywhere}.
\end{equation}
No logarithm of a potentially zero coefficient was used in this deduction.

\subsection{Recovery of the sum of antipodal densities}
For $t\ge0$, let $\ell_C(y)$ denote the length of the fiber of $C$
over $y\in u^\perp$. Convexity makes each fiber an interval, so
\[
 g_C(tu)=\int_{u^\perp}(\ell_C(y)-t)_+\,dy .
\]
Dominated convergence, and the fact that the boundary of the planar
projection has area zero, imply
\begin{equation}\label{eq:brightness}
 -\left.\frac{d}{dt}g_C(tu)\right|_{t=0+}
 =\area(C|u^\perp)
 =\frac12\int_{\Sph}|u\cdot v|\,dS_C(v).
\end{equation}
The last equality is Cauchy's projection formula. The cosine transform
is injective on even finite signed measures on the sphere; equivalently,
its restriction to the even part of a measure determines that part
\cite{Schneider}. Equality of \eqref{eq:brightness} for $K$ and $L$
therefore implies
$S_L+S_{-L}=S_K+S_{-K}$.
Using \eqref{eq:area-formula} and antipodal invariance of $\sigma$ gives
\begin{equation}\label{eq:density-sum}
 f_L(u)+f_L(-u)=f_K(u)+f_K(-u)\quad\text{almost everywhere}.
\end{equation}
The two numbers at each antipodal pair are the roots of a monic quadratic
with this sum and the product in \eqref{eq:density-product}. Hence
\eqref{eq:unordered} follows, completing the proposed proof of
Proposition~\ref{prop:geometric}.
\section{Allocation of zeros and the paired canonical product}\label{app:allocation}
Throughout this appendix $K$ is noncentrally symmetric, belongs to $C^9_+$,
and $g_L=g_K$. No differentiability of $\partial L$ is used.

\subsection{One allocation for each analytic branch}
Compactness of the sphere gives $0<w_{\min}\le w(u)\le w_{\max}<\infty$.
The real centers $\lambda_m(u)=(2m+1)\pi/w(u)$ have separation
at least $2\pi/w_{\max}$. Increase the starting index so that
the uniform error in \eqref{eq:rootsK} is smaller than a fixed fraction
of this separation. For each $u$, disjoint real-part intervals around
the $\lambda_m(u)$ then contain exactly one positive-real-part zero
of $K$ each, namely $z_m(u)$, and exclude all exceptional zeros.
Finitely many removed branches are absorbed into the bounded exceptional set.
In particular, if $z_m(u)$ is nonreal, its conjugate cannot be another
zero of $\F_K(\,\cdot\,u)$.

Choose $u_*$ with $\gamma_K(u_*)\ne0$ and increase the index again so
that every $z_m(u_*)$ is nonreal. The product
$\F_K(\zeta u_*)\F_K(-\zeta u_*)$ has a simple zero at $z_m(u_*)$:
the first factor has a simple zero and the second does not vanish,
by the preceding separation and \eqref{eq:reality}.

For each such $m$, set
\[
 P_m(u)=\F_L(z_m(u)u),\qquad Q_m(u)=\F_L(-z_m(u)u).
\]
Both functions are real analytic on the connected real-analytic manifold
$\Sph$, and \eqref{eq:product} gives $P_mQ_m=0$.
If $P_m$ is not identically zero, it is nonzero on a nonempty open set,
where $Q_m$ must vanish. The identity principle gives $Q_m\equiv0$.
Thus one factor vanishes identically. Both cannot do so, since at $u_*$
the product on the $L$ side would have a zero of multiplicity at least
two, contradicting the simple zero on the $K$ side.

Set $\varepsilon_m=1$ if $P_m\equiv0$, and $\varepsilon_m=-1$ if
$Q_m\equiv0$. In the latter case the reality symmetry implies that
$\overline{z_m(u)}$ is a zero of $\F_L(\,\cdot\,u)$.
These choices are independent of $u$.

For completeness, consider multiplicities at any direction $u$.
If $z_m(u)$ is nonreal, the product on the $K$ side has simple zeros
at $z_m(u)$ and $\overline{z_m(u)}$. The selected zero of $L$ is
therefore simple, and $L$ cannot contain the other one as well:
reality would then make both factors on the $L$ side vanish at the
selected point. If $z_m(u)$ is real, reality gives equal multiplicities
at $z_m(u)$ and $-z_m(u)$ for $L$. The $K$ product has multiplicity
two there, so each factor on the $L$ side has multiplicity one.
This handles directions where the conjugate choices coincide.

Every zero of $\F_L(\,\cdot\,u)$ is a zero of the $K$ product.
The large zeros have just been exhausted; the remaining ones lie in a
bounded set and are finite for fixed $u$, since $\F_L(0)=\vol(L)>0$.
This proves \eqref{eq:rootsL}, including the assertions about multiplicity
and exceptional zeros.

\subsection{Hadamard factorization in one direction}\label{app:hadamard}
Fix $u$. List all zeros with positive real part, including exceptions
with their multiplicities, as $\zeta_j=x_j+iy_j$. By the preceding
argument, $x_j>0$, the $y_j$ are bounded, and $x_j$ grows linearly
with $j$, after a harmless finite reindexing. The other zeros are
$-\overline{\zeta_j}$. There are no zeros on the imaginary axis by
\eqref{eq:reality}, and
\begin{equation}\label{eq:inverse-sums}
 \sum_j x_j^{-2}<\infty,\qquad
 \sum_j\left|\zeta_j^{-1}-\overline{\zeta_j}^{-1}\right|<\infty .
\end{equation}

If $L\subset B(0,T)$, then
$|F(\zeta)|\le\vol(L)e^{T|\Ima\zeta|}$, so $F$ is an entire
function of order at most one. Hadamard factorization in this case
gives a canonical product of genus at most one with an exponential
whose exponent has degree at most one \cite[p.~4]{Siu}.
Pairing its zeros and using the second sum in \eqref{eq:inverse-sums}
to absorb the linear canonical factors yields
\begin{equation}\label{eq:canonical}
 F(\zeta)=F(0)e^{iq\zeta}
 \prod_j\left(1-\frac{\zeta}{\zeta_j}\right)
        \left(1+\frac{\zeta}{\overline{\zeta_j}}\right),
 \qquad q\in\R .
\end{equation}
For clarity, the paired factor differs from one by
\[
 \frac{2iy_j\zeta-\zeta^2}{|\zeta_j|^2}.
\]
The first sum in \eqref{eq:inverse-sums} therefore proves local uniform
convergence of this product. The constant is $F(0)$ by evaluation at
zero. Both $F$ and the paired product satisfy the reality symmetry
in \eqref{eq:reality}; their quotient, an exponential of a linear
function, then has a purely imaginary coefficient. This explains
the real number $q$ in \eqref{eq:canonical}.

At $\zeta=i\tau$, each paired factor is
$[x_j^2+(\tau-y_j)^2]/|\zeta_j|^2>0$.
Thus real logarithms give
\begin{equation}\label{eq:Dseries}
 D_L(\tau)=-2q\tau+\sum_j R_{x_j,y_j}(\tau),\qquad
 R_{x,y}(\tau)=\log\frac{x^2+(\tau-y)^2}{x^2+(\tau+y)^2}.
\end{equation}
For fixed $\tau$ the logarithmic series is absolutely convergent:
its tail is $O_\tau(\sum_jx_j^{-2})$, with bounded heights $y_j$.
Taking the finite difference consequently gives the exact identity
\begin{equation}\label{eq:Hseries}
 H_L(\tau)=\sum_j\bigl(R_{x_j,y_j}(2\tau)-2R_{x_j,y_j}(\tau)\bigr).
\end{equation}
The value of $q$ is not required.

If the direction has positive finite endpoint coefficients, the
limits \eqref{eq:scaled-laplace} may be written
\[
 F(i\tau)=e^{-\alpha_L\tau}\tau^{-2}(a_L+o(1)),\qquad
 F(-i\tau)=e^{\beta_L\tau}\tau^{-2}(b_L+o(1)).
\]
Taking logarithms and then the same finite difference yields
\begin{equation}\label{eq:Hlimit}
 D_L(\tau)=-(\alpha_L+\beta_L)\tau+\log(a_L/b_L)+o(1),
 \qquad H_L(\tau)\longrightarrow\log(b_L/a_L).
\end{equation}

\section{Kernel estimates and replacement of the zeros}\label{app:kernels}
The estimates in this appendix also specify the analytic content checked,
in part, by the separate Lean files. Their use for the convex-body
transforms is supplied by the preceding and following appendices.

\subsection{A summable logarithmic remainder}
For $x\ne0$, $|y|\le M$, $\tau>0$ and $\tau\ge4M$, we have
\begin{equation}\label{eq:kernel-bound}
 \left|R_{x,y}(\tau)+\frac{4y\tau}{x^2+\tau^2}\right|
 \le\frac{36M^3\tau}{(x^2+\tau^2)^2}
 \le\frac{36M^3}{\tau x^2}.
\end{equation}
Here is a direct proof. Put $A=x^2+\tau^2$, $B=A+y^2$ and
$v=2\tau y/B$. Then $|v|\le2|y|/\tau\le1/2$, and
\[
 R_{x,y}(\tau)=\log(1-v)-\log(1+v).
\]
The derivative of $\log(1-v)-\log(1+v)+2v$ is
$-2v^2/(1-v^2)$. Integration from zero therefore gives, for
$|v|\le1/2$,
\[
 |\log(1-v)-\log(1+v)+2v|\le4|v|^3.
\]
Since $\tau^2\le A$ and $B\ge A$, this bound is at most
$32M^3\tau/A^2$. The denominator correction is
\[
 \left|\frac{4\tau y}{A}-\frac{4\tau y}{B}\right|
 =\frac{4\tau|y|^3}{AB}\le\frac{4M^3\tau}{A^2}.
\]
Their sum proves the first inequality in \eqref{eq:kernel-bound};
the second follows from $A^2\ge\tau^2x^2$.

For the zeros in Appendix~\ref{app:hadamard}, summing the bound proves
absolute summability of the remainders and
\begin{equation}\label{eq:linearize}
 \sum_jR_{x_j,y_j}(\tau)
 =-4\sum_j\frac{y_j\tau}{x_j^2+\tau^2}+o(1).
\end{equation}
The error is bounded in modulus by
$36M^3\tau^{-1}\sum_jx_j^{-2}$.
The rational series itself is absolutely convergent, since its summands
are bounded in modulus by $M\tau/x_j^2$.
Apply \eqref{eq:linearize} at $\tau$ and $2\tau$ in
\eqref{eq:Hseries}. The algebraic identity
\begin{equation}\label{eq:difference-kernel}
 -\frac{8y\tau}{x^2+4\tau^2}
 +\frac{8y\tau}{x^2+\tau^2}
 =\frac{24y\tau^3}{(x^2+\tau^2)(x^2+4\tau^2)}
\end{equation}
gives \eqref{eq:Hkernel}. All estimates used here are for sums of
absolute values; no exchange of conditionally convergent series is involved.

\subsection{Positive weights and bounded mass}\label{app:weights}
Introduce
\[
 k_\tau(x)=\frac{\tau}{x^2+\tau^2},\qquad
 \eta_\tau(x)=\frac{24\tau^3}{(x^2+\tau^2)(x^2+4\tau^2)}
            =8k_\tau(x)-4k_{2\tau}(x).
\]
For $\tau>0$ and $x>0$,
\begin{equation}\label{eq:positive-weight}
 0<\eta_\tau(x)\le8k_\tau(x)\le\frac8\tau .
\end{equation}
If $x_m\ge\rho m$ for $m\ge m_1\ge1$ and $\rho>0$, monotonicity
of $t\mapsto\tau/(\rho^2t^2+\tau^2)$ gives
\[
 \sum_{m\ge m_1}k_\tau(x_m)
 \le\sum_{m\ge1}\frac{\tau}{\rho^2m^2+\tau^2}
 \le\int_0^\infty\frac{\tau\,dt}{\rho^2t^2+\tau^2}
 =\frac{\pi}{2\rho}.
\]
Finitely many initial terms are bounded by their number divided by
$\tau$. Thus the total mass of both $k_\tau(x_m)$ and
$\eta_\tau(x_m)$ is bounded for $\tau\ge1$.
Every individual weight tends to zero.

It follows that if $e_m\to0$, then
\begin{equation}\label{eq:height-error}
 \sum_m e_m\eta_\tau(x_m)\longrightarrow0 .
\end{equation}
Indeed $e_m$ is bounded, so the series converges absolutely for each
$\tau$. Given $\epsilon>0$, choose an index after which $|e_m|<\epsilon$.
The tail has modulus at most $\epsilon$ times the uniform mass bound;
the finite head tends to zero. Let $\epsilon\downarrow0$.
This argument is only directional: its constants may depend on the
fixed $u$, which is sufficient for the subsequent common-parameter argument.

\subsection{Replacing the real parts without differentiating a series}
For $x,a,\tau>0$ with $|x-a|\le D$,
\begin{equation}\label{eq:real-replacement}
 |k_\tau(x)-k_\tau(a)|
 \le D\left(\frac1{x^2}+\frac1{a^2}\right).
\end{equation}
To see this, write $A=x^2+\tau^2$, $B=a^2+\tau^2$. The difference
has modulus
$|x-a|\tau(x+a)/(AB)$, and $\tau(x+a)\le A+B$.
This gives $D(1/A+1/B)$ and then \eqref{eq:real-replacement}.
For fixed $x,a$ the difference tends to zero as $\tau\to\infty$.

If $(x_m)$ and $(a_m)$ are positive, their difference is uniformly
bounded, and $\sum_m(x_m^{-2}+a_m^{-2})<\infty$, dominated convergence
for series gives
\[
 \sum_m|k_\tau(x_m)-k_\tau(a_m)|\longrightarrow0.
\]
The same holds for $\eta_\tau$, by its expression in terms of
$k_\tau$ and $k_{2\tau}$; bounded signed coefficients may also be
inserted.

Apply these facts to a fixed direction and
$a_m=\lambda_m=(2m+1)\pi/w$. The assumptions follow from
\eqref{eq:rootsL}; in particular $x_m$ and $\lambda_m$ grow linearly.
First replace $y_m$ by $\varepsilon_m\gamma_K(u)$ using
\eqref{eq:height-error}. Then replace $x_m$ by $\lambda_m$ using
\eqref{eq:real-replacement}. The exceptional zeros are finite and
each contributes $o(1)$ by \eqref{eq:positive-weight}.
Consequently \eqref{eq:Hkernel} becomes
\begin{equation}\label{eq:model-sum}
 H_L(\tau)=\gamma_K(u)\sum_{m\ge m_0}
                         \varepsilon_m\eta_\tau(\lambda_m)+o(1).
\end{equation}

\section{The common mean and proof of the transport proposition}\label{app:mean}
We finish the proposed proof of Proposition~\ref{prop:transport}.
For fixed $u$, put $s=2\pi/w(u)$, so that
$\lambda_m=s(m+\tfrac12)$ and $\tau=sR$. A direct calculation gives
\[
 \eta_{sR}\bigl(s(m+\tfrac12)\bigr)
 =\frac{12w(u)}{\pi}
   \frac{R^3}{((m+\tfrac12)^2+R^2)((m+\tfrac12)^2+4R^2)} .
\]
Thus \eqref{eq:model-sum} gives precisely \eqref{eq:Hmean}, with
$N(R)$ defined in \eqref{eq:mean}. Its only input from the direction
is the common sequence of signs, which Appendix~\ref{app:allocation}
has already shown to be independent of that direction.

\subsection{Normalization of the positive weights}
Let
\[
 \varphi(t)=\frac1{(1+t^2)(4+t^2)},\qquad t\ge0.
\]
The mass of the weights of $N$ is
\begin{equation}\label{eq:mass}
 W(R)=\frac{12}{\pi R}
           \sum_{m\ge m_0}\varphi((m+\tfrac12)/R).
\end{equation}
The function $\varphi$ is nonnegative, decreasing and integrable.
Left and right rectangular sums on intervals of length $1/R$
bound its integral. The midpoint sum is between the same two sums,
whose difference is $\varphi(0)/R$. Therefore
\[
 \left|\frac1R\sum_{m\ge0}\varphi((m+\tfrac12)/R)
             -\int_0^\infty\varphi(t)\,dt\right|
 \le\frac{\varphi(0)}R .
\]
Removing the first $m_0$ terms changes the expression by at most
$m_0\varphi(0)/R$. Partial fractions give
\[
 \int_0^\infty\varphi(t)\,dt
 =\frac13\int_0^\infty
       \left(\frac1{1+t^2}-\frac1{4+t^2}\right)\,dt
 =\frac{\pi}{12}.
\]
It follows that $W(R)\to1$. In particular,
\begin{equation}\label{eq:signed-bound}
 |N(R)|\le W(R),\qquad \limsup_{R\to\infty}|N(R)|\le1 .
\end{equation}

\subsection{Existence and independence of the limit}
Let $G$ be a full-measure set where the positive endpoint expansions hold.
There is $u_0\in G$ with $\gamma_K(u_0)\ne0$: a nonempty open subset of
the nonzero set of the continuous function $\gamma_K$ has positive
surface measure. Equations \eqref{eq:Hlimit} and \eqref{eq:Hmean} then give
\[
 N(R)=\frac{H_L(2\pi R/w(u_0))}{w(u_0)\gamma_K(u_0)}+o(1)
 \longrightarrow
 c:=\frac{\log(b_L(u_0)/a_L(u_0))}
          {w(u_0)\gamma_K(u_0)} .
\]
By \eqref{eq:signed-bound}, $c\in[-1,1]$.
For any other $u\in G$, the very same $N(R)$ appears in
\eqref{eq:Hmean}, so
\[
 \log(b_L(u)/a_L(u))=w(u)\gamma_K(u)c.
\]
This includes $\gamma_K(u)=0$, since then the right side of
\eqref{eq:Hmean} tends to zero. The parameter is unique because the
identity holds at $u_0$. This is \eqref{eq:transport}.

The common mean, rather than pointwise convergence of individual signs
$\varepsilon_m$, is the quantity used in the proof. Its existence is forced
by the endpoint limit in a single admissible direction. Neither convergence
of the sign sequence nor uniform endpoint asymptotics for $L$ is assumed.
\begin{thebibliography}{9}
\bibitem{Bianchi2016}
G.~Bianchi, \emph{The covariogram and Fourier--Laplace transform in
$\mathbb C^n$}, Proceedings of the London Mathematical Society
113 (2016), 1--23.
\href{https://doi.org/10.1112/plms/pdw020}{doi:10.1112/plms/pdw020};
\href{https://arxiv.org/abs/1312.7816}{arXiv:1312.7816}.

\bibitem{Bianchi2023}
G.~Bianchi, \emph{The covariogram problem}, in
\emph{Harmonic Analysis and Convexity} (2023), 37--82.
\href{https://doi.org/10.1515/9783110775389-002}{doi:10.1515/9783110775389-002};
\href{https://people.dimai.unifi.it/bianchi/lavori/covariogram_survey14.pdf}{author's manuscript}.

\bibitem{Schneider}
R.~Schneider, \emph{Convex Bodies: The Brunn--Minkowski Theory},
second expanded edition, Encyclopedia of Mathematics and its Applications
151, Cambridge University Press, 2014.
\href{https://doi.org/10.1017/CBO9781139003858}{doi:10.1017/CBO9781139003858}.

\bibitem{EvansGariepy}
L.~C.~Evans and R.~F.~Gariepy, \emph{Measure Theory and Fine Properties
of Functions}, second edition, CRC Press, 2025.
\href{https://www.routledge.com/Measure-Theory-and-Fine-Properties-of-Functions/Evans/p/book/9781032946443}{publisher's description and contents}.

\bibitem{Siu}
Y.-T.~Siu, \emph{Weierstrass and Hadamard Factorization of Entire Functions},
Math 213a course notes, Harvard University, Fall 2024.
\href{https://people.math.harvard.edu/~siu/math213a/weierstrass_and_hadamard_factorization_of_entire_function.pdf}{author's notes}.
\end{thebibliography}
\end{document}
